% Generated from locale/tr/content/turing-machines/machines-computations/representing-tms.tex; do not hand-edit.


\olfileid[tr]{tur}{mac}{rep}
\olsection{Turing Makinelerinin Gösterilmesi}

\begin{explain}
Turing makineleri \emph{durum diyagramları}yla görsel olarak gösterilebilir.
Diyagramlar oklarla bağlanan durum hücrelerinden oluşur. Bekleneceği üzere
her durum hücresi makinenin bir durumunu gösterir. Her ok, o durumda yerine
getirilebilecek bir yönergeyi; okun üstüne ya da altına yazılan bilgi ise
yönergenin ayrıntılarını gösterir. Yalnızca iki iç durumu, $q_0$ ve $q_1$,
ve tek yönergesi olan aşağıdaki makineyi ele alalım:
\[
\begin{tikzpicture}[->,>=stealth',shorten >=1pt,auto,node distance=2.8cm,
                    semithick]
  \tikzstyle{every state}=[fill=none,draw=black,text=black]

  \node[initial,state] (A) {$q_0$};
  \node[state] (B) [right of=A] {$q_1$};

  \path (A) edge node {\TMtrans{\TMblank}{\TMstroke}{\TMright}} (B);
\end{tikzpicture}
\]
Turing makinesinin bir okuma-yazma kafasına ve üzerinde girdi yazılı bir
şeride sahip olduğunu anımsayın. Yönerge şöyle okunabilir:
\emph{$q_0$ durumunda $\TMblank$ okunuyorsa $\TMstroke$ yaz, sağa git ve
$q_1$ durumuna geç}. Bu, $\tuple{q_0,\TMblank}$ çiftini
$\tuple{q_1,\TMstroke,\TMright}$ üçlüsüne götüren geçiş fonksiyonuna
eşdeğerdir.
\end{explain}

\begin{ex}
\emph{Çift Sayı Makinesi}: Aşağıdaki Turing makinesi şeritte ancak ve ancak
çift sayıda $\TMstroke$ varsa durur; burada bütün $\TMstroke$ simgelerinin
şeritteki ilk $\TMblank$ simgesinden önce geldiğini varsayıyoruz.
\[
\begin{tikzpicture}[->,>=stealth',shorten >=1pt,auto,node distance=2.8cm,
                    semithick]
  \tikzstyle{every state}=[fill=none,draw=black,text=black]

  \node[initial,state] (A) {$q_0$};
  \node[state] (B) [right of=A] {$q_1$};

  \path (A) edge [bend left] node {\TMtrans{\TMstroke}{\TMstroke}{\TMright}} (B)
        (B) edge [loop above] node {\TMtrans{\TMblank}{\TMblank}{\TMright}} (B)
            edge [bend left] node {\TMtrans{\TMstroke}{\TMstroke}{\TMright}} (A);
\end{tikzpicture}
\]

Durum diyagramı aşağıdaki geçiş fonksiyonuna karşılık gelir:
\begin{align*}
\delta(q_0,\TMstroke) & = \tuple{q_1,\TMstroke,\TMright},\\
\delta(q_1,\TMstroke) & = \tuple{q_0,\TMstroke,\TMright},\\
\delta(q_1,\TMblank) & = \tuple{q_1,\TMblank,\TMright}.
\end{align*}
\end{ex}

\begin{explain}
Yukarıdaki makine yalnızca girdi çift sayıda çizgiden oluştuğunda durur.
Aksi durumda makine kuramsal olarak sonsuza kadar çalışmayı sürdürür. Her
makine ve girdi için çıktıyı belirlemek üzere makinenin
\emph{konfigürasyonlarını} adım adım izlemek mümkündür. Konfigürasyonların
biçimsel tanımını daha sonra vereceğiz. Şimdilik onları, çalışma sırasında
herhangi bir andaki makine durumunu gösteren bir diyagram dizisi olarak
sezgisel biçimde düşünebiliriz. Konfigürasyonlar şerit içeriğini, makinenin
durumunu ve okuma-yazma kafasının konumunu gösterir.

Dört $\TMstroke$ girdisiyle başlatılan çift sayı makinesinin
konfigürasyonlarını izleyelim. Bu durumda makinenin durmasını bekleriz.
Ardından makineyi üç $\TMstroke$ girdisiyle çalıştıracağız; o zaman makine
sonsuza kadar çalışacaktır.

Makine $q_0$ durumunda, en soldaki $\TMstroke$ simgesini tarayarak başlar.
Makinenin başlangıç durumunu şöyle gösterebiliriz:
\[
\TMendtape \TMstroke_0 \TMstroke \TMstroke \TMstroke \TMblank \ldots
\]
Yukarıdaki konfigürasyon açıktır. Makine birinci durumda, en soldaki
$\TMstroke$ simgesini tarayarak başlar. Bu, ilk $\TMstroke$ üzerindeki durum
adının alt indisiyle gösterilir. Bu noktada uygulanabilir yönerge
$\delta(q_0,\TMstroke)=\tuple{q_1,\TMstroke,\TMright}$ olduğundan makine
şeritte sağa gider ve $q_1$ durumuna geçer.
\[
\TMendtape \TMstroke \TMstroke_1 \TMstroke \TMstroke \TMblank \ldots
\]
Makine artık $q_1$ durumunda $\TMstroke$ taradığından
$\delta(q_1,\TMstroke)=\tuple{q_0,\TMstroke,\TMright}$ yönergesini
``izlemeliyiz''. Bunun sonucunda
\[
\TMendtape \TMstroke \TMstroke \TMstroke_0 \TMstroke \TMblank \ldots
\]
konfigürasyonu elde edilir. Makine çalışmayı sürdürürken kurallar aynı sırayla
yeniden uygulanır ve aşağıdaki iki konfigürasyon ortaya çıkar:
\[
\TMendtape \TMstroke \TMstroke \TMstroke \TMstroke_1 \TMblank \ldots
\]
\[
\TMendtape \TMstroke \TMstroke \TMstroke \TMstroke \TMblank_0 \ldots
\]
Makine artık $q_0$ durumunda $\TMblank$ taramaktadır. Geçiş diyagramından
yerine getirilecek bir yönerge olmadığı kolayca görülür; dolayısıyla makine
durmuştur. Bu, girdinin kabul edildiği anlamına gelir.

Şimdi makineyi üç $\TMstroke$ girdisiyle başlattığımızı varsayalım. Aynı
yönergeler yerine getirildiğinden ilk birkaç konfigürasyon, yalnızca şerit
girdisindeki küçük farkla benzerdir:
\[
\TMendtape \TMstroke_0 \TMstroke \TMstroke \TMblank \ldots
\]
\[
\TMendtape \TMstroke \TMstroke_1 \TMstroke \TMblank \ldots
\]
\[
\TMendtape \TMstroke \TMstroke \TMstroke_0 \TMblank \ldots
\]
\[
\TMendtape \TMstroke \TMstroke \TMstroke \TMblank_1 \ldots
\]
Makine bütün $\TMstroke$ simgelerini geçmiş ve $q_1$ durumunda bir
$\TMblank$ okumaktadır. Diyagramda gösterildiği gibi
$\delta(q_1,\TMblank)=\tuple{q_1,\TMblank,\TMright}$ biçiminde bir yönerge
vardır. Şerit sağa doğru sonsuza kadar $\TMblank$ ile dolu olduğundan makine
bu yönergeyi \emph{sonsuza kadar} yerine getirir; $q_1$ durumunda kalıp
giderek daha sağa gider. Makine hiçbir zaman durmaz ve girdiyi kabul etmez.
\end{explain}

\begin{explain}
Bütün makinelerin durmayacağını belirtmek önemlidir. Durma, makinenin yerine
getirecek yönergesinin kalmaması demekse her durumdaki her simge için çıkan
bir ok bulunmasını güvenceye alarak hiç durmayan bir makine kurabiliriz.
Çift sayı makinesi, $q_0$ durumunda $\TMblank$ taramak için bir yönerge
eklenerek sonsuza kadar çalışacak biçimde değiştirilebilir.
\end{explain}

\begin{ex}
\[
\begin{tikzpicture}[->,>=stealth',shorten >=1pt,auto,node distance=2.8cm,
                    semithick]
  \tikzstyle{every state}=[fill=none,draw=black,text=black]

  \node[initial,state] (A) {$q_0$};
  \node[state] (B) [right of=A] {$q_1$};

  \path
  (A) edge [bend left] node {\TMtrans{\TMstroke}{\TMstroke}{\TMright}} (B)
      edge [loop above] node {\TMtrans{\TMblank}{\TMblank}{\TMright}} (A)
  (B) edge [loop above] node {\TMtrans{\TMblank}{\TMblank}{\TMright}} (B)
      edge [bend left] node {\TMtrans{\TMstroke}{\TMstroke}{\TMright}} (A);
\end{tikzpicture}
\]
\end{ex}

\begin{explain}
Makine tabloları Turing makinelerini göstermenin başka bir yoludur. Makine
tablolarında şerit alfabesi $x$ ekseninde, makine durumları kümesi ise $y$
ekseninde gösterilir. Tabloda her durumla simgenin kesişimine yönergenin
geri kalanı, yani yeni durum, yeni simge ve hareket yönü yazılır. Makine
tabloları makinenin hangi durumda ve hangi simgede durduğunu belirlemeyi
kolaylaştırır. Tablodaki her boşluk makinenin olası bir durma noktasıdır.
Makinenin durduğu noktaların her zaman hemen görünmediği durum diyagramları
ve yönerge kümelerinden farklı olarak, durma noktaları makine tablosundaki
boşluklar bulunarak hızla belirlenir.
\end{explain}

\begin{ex}
Çift sayı makinesinin makine tablosu şöyledir:
\[
\centering
\begin{tabular}{lllll}
\cline{2-4}
\multicolumn{1}{l|}{} & \multicolumn{1}{c|}{$\TMblank$}
& \multicolumn{1}{c|}{$\TMstroke$} & \multicolumn{1}{c|}{$\TMendtape$} & \\ \cline{1-4}
\multicolumn{1}{|l|}{$q_0$} & \multicolumn{1}{l|}{}
& \multicolumn{1}{l|}{\TMtrans{\TMstroke}{q_1}{\TMright}} &
\multicolumn{1}{l|}{\phantom{\TMtrans{\TMstroke}{q_1}{\TMright}}} & \\ \cline{1-4}
\multicolumn{1}{|l|}{$q_1$} & \multicolumn{1}{l|}{\TMtrans{\TMblank}{q_1}{\TMright}}
& \multicolumn{1}{l|}{\TMtrans{\TMstroke}{q_0}{\TMright}} &
\multicolumn{1}{l|}{\phantom{\TMtrans{\TMstroke}{q_1}{\TMright}}} & \\ \cline{1-4}
\end{tabular}
\]
Görüldüğü gibi makine $q_0$ durumunda $\TMblank$ tararken durur.
\end{ex}

\begin{explain}
Şimdiye kadar yalnızca girdiyi okuyup kabul eden makineleri ele aldık. Fakat
Turing makineleri hem okuma hem yazma yeteneğine sahiptir. Böyle bir makine
örneği, çok sayıda başka örnek de vardır, bir \emph{iki katına çıkarıcı}dır.
İki katına çıkarıcı şeritte $n$ tane $\TMstroke$ bloğuyla başlatıldığında
$2n$ tane $\TMstroke$ bloğunu çıktı olarak verir.
\end{explain}

\begin{ex}
  \ollabel{ex:doubler}
  İki katına çıkarıcı makineyi kurmadan önce problemi çözecek bir
  \emph{strateji} bulmak önemlidir. Makine, bizim kurduğumuz biçimiyle kaç
  tane $\TMstroke$ okuduğunu anımsayamadığından şeritteki bütün
  $\TMstroke$ simgelerini izlemenin bir yoluna gereksinim duyarız. Böyle bir
  yol, çıktıyı girdiden bir $\TMblank$ ile ayırmaktır. Makine girdideki ilk
  $\TMstroke$ simgesini silebilir, girdinin geri kalanını geçebilir, bir
  $\TMblank$ bırakıp iki yeni $\TMstroke$ yazabilir. Ardından geri dönüp
  girdideki ikinci $\TMstroke$ simgesini bulur ve onu da iki katına çıkarır.
  Her bir girdi $\TMstroke$ simgesi için iki çıktı $\TMstroke$ simgesi yazar.
  Girdiyi ilerledikçe silerek hiçbir $\TMstroke$ simgesinin atlanmamasını ya
  da iki kez iki katına çıkarılmamasını güvenceye alırız. Bütün girdi
  silindiğinde şeritte $2n$ tane $\TMstroke$ kalır. Ortaya çıkan Turing
  makinesinin durum diyagramı \olref{fig:doubler} içinde gösterilmiştir.
  \begin{figure}
\[
\begin{tikzpicture}[->,>=stealth',shorten >=1pt,auto,node distance=2.8cm,
                    semithick]
  \tikzstyle{every state}=[fill=none,draw=black,text=black]

  \node[initial,state] (1) {$q_0$};
  \node[state] (2) [right of=1] {$q_1$};
  \node[state] (3) [right of=2] {$q_2$};
  \node[state] (4) [below of=3] {$q_3$};
  \node[state] (5) [left of=4] {$q_4$};
  \node[state] (6) [left of=5] {$q_5$};

  \path (1) edge node {\TMtrans{\TMstroke}{\TMblank}{\TMright}} (2)
    (2) edge [loop above] node {\TMtrans{\TMstroke}{\TMstroke}{\TMright}} (2)
      edge node {\TMtrans{\TMblank}{\TMblank}{\TMright}} (3)
    (3) edge [loop above] node {\TMtrans{\TMstroke}{\TMstroke}{\TMright}} (3)
        edge node {\TMtrans{\TMblank}{\TMstroke}{\TMright}} (4)
    (4) edge [loop below] node {\TMtrans{\TMblank}{\TMstroke}{\TMleft}} (4)
        edge node {\TMtrans{\TMstroke}{\TMstroke}{\TMleft}} (5)
    (5) edge [loop below] node {\TMtrans{\TMstroke}{\TMstroke}{\TMleft}} (5)
        edge node {\TMtrans{\TMblank}{\TMblank}{\TMleft}} (6)
    (6) edge [loop below] node {\TMtrans{\TMstroke}{\TMstroke}{\TMleft}} (6)
        edge node {\TMtrans{\TMblank}{\TMblank}{\TMright}} (1);
\end{tikzpicture}
\]
\caption{İki katına çıkarıcı makine}
\ollabel{fig:doubler}
\end{figure}
\end{ex}

\begin{prob}
Keyfî bir girdi seçip \olref[tur][mac][rep]{ex:doubler} içindeki iki katına
çıkarıcı makinenin konfigürasyonlarını adım adım izleyin.
\end{prob}

\begin{prob}
Alfabesi $\{\TMendtape,\TMblank,A,B\}$ olan ve $A$ sayısıyla $B$ sayısının
eşit \emph{ve} bütün $A$ simgelerinin bütün $B$ simgelerinden önce geldiği
her $A$ ve $B$ dizgesini kabul eden, yani üzerinde duran; $A$ ile $B$
sayılarının eşit olmadığı ya da bütün $A$ simgelerinin bütün $B$
simgelerinden önce gelmediği her dizgeyi reddeden, yani üzerinde durmayan bir
Turing makinesi tasarlayın. Örneğin makine $AABB$ ile $AAABBB$ girdilerini
kabul, $AAB$ ile $AABBAABB$ girdilerini reddetmelidir.
\end{prob}

\begin{prob}
Alfabesi $\{\TMendtape,\TMblank,A,B\}$ olan, girdi olarak herhangi bir
$A$ ve $B$ dizgesi $\alpha$ alıp $\alpha\alpha$ biçiminde çıktı üretmek için
onu çoğaltan bir Turing makinesi tasarlayın. Örneğin $ABBA$ girdisi
$ABBAABBA$ çıktısını vermelidir.
\end{prob}

\begin{prob}
\emph{Alfabetik mi?:} Alfabesi $\{\TMendtape,\TMblank,A,B\}$ olan ve sonlu
bir $A$ ve $B$ dizisi girdi olarak verildiğinde bütün $A$ simgelerinin bütün
$B$ simgelerinin solunda olup olmadığını denetleyen bir Turing makinesi
tasarlayın. Makine girdi dizgesini şeritte bırakmalı; dizge ``alfabetik'' ise
durmalı, değilse sonsuza kadar döngüye girmelidir.
\end{prob}

\begin{prob}
\emph{Alfabetik sıralayıcı:} Alfabesi
$\{\TMendtape,\TMblank,A,B\}$ olan ve sonlu bir $A$ ve $B$ dizisini girdi
olarak alıp bütün $A$ simgeleri bütün $B$ simgelerinin soluna gelecek biçimde
yeniden düzenleyen bir Turing makinesi tasarlayın. Örneğin $BABAA$ dizisi
$AAABB$, $ABBABB$ dizisi ise $AABBBB$ olmalıdır.
\end{prob}

